\subsection{The same theta input and the cubic-fourfold discriminant}

The preceding theta construction also has an exact interpretation for
ordinary cubic fourfolds. The character changes the answer: on the
orthogonal domain the weight-$48$ form must be squared before applying
the invariant-ring isomorphism. This is a fixed-dimensional
compatibility of known constructions, with an explicit discriminant
identity; no historical priority is claimed.

\paragraph{Primary geometric inputs and conventions.}
Let $V=\mathbb C^6$, $\mathscr S_6=\operatorname{Sym}^3V^*$,
and write $\Delta_6$ for its discriminant. Put
\[
 M=U^{\oplus2}\oplus E_8(-1)^{\oplus2}\oplus A_2(-1),
 \qquad K=E_6.
\]
Thus $M$ has signature $(2,20)$ and embeds primitively in
$L=II_{2,26}$ with complement $K(-1)$, exactly as in P-F46.
Let $D^+$ be one connected component of its type-IV domain and
$\Gamma^+=\widetilde O^+(M)$ the group preserving that component
and acting trivially on $A_M=M^\vee/M$.

Looijenga, \emph{The period map for cubic fourfolds}, Section~1
printed p.~3 and Section~2 p.~7, identifies the group for the
cubic-fourfold period map with this discriminant kernel, allowing
both components before restriction to $D^+$. Theorem~3.1,
printed p.~11, gives an isomorphism of algebras
\[
 \mathcal P_6:\bigoplus_{k\ge0}
 H^0(\dot D^+,\mathcal A(k))^{\Gamma^+}
       \ \xrightarrow{\ \sim\ }
 \mathbb C[\mathscr S_6]^{SL(V)},
 \qquad \deg\mathcal P_6(h)=2\operatorname{wt}(h).
 \tag{P-F55a}
\]
Here $\dot D^+$ and its section spaces are precisely the
arrangement-complement spaces in that theorem. In particular this
notation retains its automorphic section convention. The theorem
identifies the line bundles over the locus of cubics with at most
ADE singularities. Laza, \emph{The moduli space of cubic fourfolds
via the period map}, Theorem~1.1, printed p.~1, identifies this
image with the complement of $\mathcal H_\infty$, and the smooth
image with the complement of
$\mathcal H_\infty\cup\mathcal H_\Delta$.
The source URLs are
\url{https://arxiv.org/pdf/0705.0951} and
\url{https://arxiv.org/pdf/0705.0949}.

Let $g=\Theta_K/\delta$ denote the vector-valued form, with the
discriminant-coset identification and Weil convention of (P-F46a).
Write $\Psi=\Psi_M(g)$. By the exact theta contraction of P-F46,
this is, up to a nonzero scalar, the $E_6$ quasi-pullback of
$\Phi_{12}$. Laza, Notation~6.4 printed p.~24 and
Proposition~6.6 p.~25, uses the same primitive embedding and
quasi-pullback: its weight is $48$ and its orders on
$\mathcal H_\Delta,\mathcal H_\infty$ are $1,27$ respectively.

\begin{corollary}[P-F55: an exact second invariant image and its cyclic compatibility]
\label{cor:fourfold-theta-discriminant}
The multiplier of $\Psi$ on $\Gamma^+$ is $\det$, and
\[
          \mathcal P_6(\Psi^2)=c_6\Delta_6,
                    \qquad c_6\in\mathbb C^*.
 \tag{P-F55b}
\]
The square has automorphic weight $96$, ordinary coefficient
degree $192$, and canonical stretching degree $96$ in the
actual cubic plethysm algebra
\[
       T^{(6)}=\bigoplus_{s\ge0}
        \left(\operatorname{Sym}^{2s}
                    (\operatorname{Sym}^3\mathbb C^6)\right)^{SL_6}.
 \tag{P-F55c}
\]
Its $GL_6$ change-of-variables character is $\det^{96}$.
The same identification respects all products in
$\mathbb C[\Psi^2]\cong\mathbb C[\Delta_6]$.

For a cubic $f$ in five variables and an additional variable $w$,
there is a nonzero constant $C$, depending only on the two
discriminant normalizations, such that the following polynomial
identity holds in $a$ and all coefficients of $f$:
\[
           \Delta_6(a w^3-f)=C a^{32}\Delta_5(f)^2.
 \tag{P-F55d}
\]
With the discriminant normalization of Bus\'e--Jouanolou
Definition~4.6, the constant is $C=3^{33}$.
At $a=1$ this agrees, up to scalar normalization, with restriction
of the named orthogonal form to the cubic-threefold ball:
\[
        \Psi^2|_{\mathbb B}\doteq\phi^2,
        \qquad \mathcal P_5(\phi^2)\doteq\Delta_5^2.
 \tag{P-F55e}
\]
Here $\phi$ and $\mathcal P_5$ are those of P-F41, and
$\doteq$ denotes equality up to a nonzero scalar. This last
statement concerns the named forms and their powers; it does not
assert that the affine substitution $f\mapsto w^3-f$ is a graded
homomorphism between the full invariant rings.
\end{corollary}

\begin{proof}
We first check the character on the exact group of (P-F55a).
Gritsenko, \emph{24 faces of the Borcherds modular form
$\Phi_{12}$}, printed p.~2, states
\[
             \Phi_{12}\in M_{12}(O^+(II_{2,26}),\det).
\]
The source is \url{https://arxiv.org/pdf/1203.6503}.
The gluing subgroup for
$M\oplus K(-1)\subset L$ is the graph of the discriminant
isometry used in P-F46. For $\gamma\in\Gamma^+$,
$\gamma\oplus\mathrm{id}_K$ acts trivially on this graph, hence
extends to an isometry $\widetilde\gamma$ of $L$.
It preserves the chosen domain component, since it preserves
the positive-plane component containing $D^+$, and
$\det\widetilde\gamma=\det\gamma$.
In the quasi-pullback one divides $\Phi_{12}$ by the product
of one linear form for each pair of opposite roots in $K(-1)$,
and then restricts to $D^+$. There are $72/2=36$ such factors.
The extension fixes each of these linear forms, not merely
their collection, because its restriction to $K(-1)$ is the
identity. Thus the quotient and its nonzero restriction have
character $\det\gamma$. This proves the multiplier assertion
and the trivial character of $\Psi^2$.

For clarity, passing to one domain component changes no invariant
section ring in Looijenga's theorem. His $D$ has two components
interchanged by an element of the discriminant kernel
(Section~2, p.~7). A section invariant under the stabilizer
of $D^+$ extends uniquely to the other component by such an
element; independence of the chosen element follows from
invariance under that stabilizer. This identifies the ring in
the cited theorem with the ring in (P-F55a).

The form $\Psi^2$ is holomorphic on the full orthogonal domain,
so is an allowed section in the arrangement-complement ring.
The isomorphism (P-F55a) sends it to a nonzero homogeneous
invariant $I$ of coefficient degree $2\cdot96=192$.
We next identify this polynomial without assuming a divisor
multiplicity convention on an orbifold quotient.

The discriminant in $\mathscr S_6$, whose dimension is $56$,
has a dense locus of cubics with a unique ordinary node.
Here is a direct verification of the density. The incidence
\[
  \mathscr I=\{(F,p)\in\mathscr S_6\times\mathbb P^5:
                                  dF(p)=0\}
\]
is a rank-$50$ vector bundle over $\mathbb P^5$. At
$p=[1:0:\cdots:0]$ its six independent conditions remove
the coefficients of $x_0^3$ and $x_0^2x_i$, $1\le i\le5$.
It is therefore irreducible of dimension $55$, and projection
to $\mathscr S_6$ is proper. Requiring two distinct singular
points imposes twelve independent conditions: after moving
them to the first two coordinate points, the two lists of
six degree-three monomials are disjoint. The corresponding
ordered-pair incidence has dimension $10+(56-12)=54$.
The condition that the quadratic part at the marked point be
degenerate is a proper closed condition on $\mathscr I$.
Properness follows, for example, from
\[
        F_0=x_0\sum_{i=1}^5x_i^2+\sum_{i=1}^5x_i^3.
 \tag{P-F55f}
\]
If $x_0=0$, its gradient equations force all $x_i=0$.
Otherwise set $x_0=1$. Each $x_i$ is then $0$ or $-2/3$,
and $\sum_i x_i^2=0$ forces all to be zero. Hence its only
singular point is $[1:0:\cdots:0]$, with nondegenerate affine
quadratic part $\sum_i x_i^2$.
The degenerate-quadratic locus and the closure of the
two-point locus have images of dimension at most $54$.
Their complement in the discriminant is the asserted dense
ordinary-node locus. In particular the incidence image has
dimension $55$, also following from the zero-dimensional
fiber just exhibited.

These nodal cubics have ADE singularities and are stable by
Looijenga's Theorem~3.2, p.~11. Their periods lie on
$\mathcal H_\Delta$ and outside $\mathcal H_\infty$ by
the smooth and ADE image descriptions above.
Since $\Psi$ vanishes on $\mathcal H_\Delta$, the
line-bundle identification in Theorem~3.1 shows that $I$
vanishes on this dense nodal locus.

Bus\'e--Jouanolou, \emph{A computational approach to the
discriminant of homogeneous polynomials}, Definition~4.6 and
Proposition~4.7(i), printed p.~40, and Theorem~4.26, p.~53,
show that over $\mathbb C$ the cubic discriminant is a prime
polynomial cutting out this hypersurface, with coefficient
degree $6(3-1)^5=192$; see
\url{https://arxiv.org/pdf/1210.4697}.
Thus $\Delta_6$ divides $I$ in the polynomial coefficient
ring. Their degrees agree, so $I=c_6\Delta_6$, with
$c_6\ne0$. This proves (P-F55b).
The same source, Proposition~4.13, pp.~42--43, gives
\[
          \Delta_6(F\circ h)=\det(h)^{3\cdot2^5}\Delta_6(F),
                       \qquad h\in GL_6,
 \tag{P-F55g}
\]
which is the claimed character.
The center of $SL_6$ permits only even invariant coefficient
degrees, and the canonical regrading is $d_0=6/\gcd(6,3)=2$.
The actual Cartan identification in P-F19 therefore gives
(P-F55c) with degree $192/2=96$. The algebra isomorphism
respects multiplication. Since a nonconstant polynomial is
transcendental over $\mathbb C$ inside a polynomial ring,
the subalgebras generated by the two named elements are
polynomial algebras and are identified as asserted.

The cyclic discriminant formula is already known, and can
be applied directly. Bus\'e--Jouanolou, Example~4.10,
printed p.~41, gives in their normalization
\[
 \operatorname{Disc}(a w^d+h)
   =d^{(d-1)^{n-1}+(-1)^n}
       a^{(d-1)^{n-1}}\operatorname{Disc}(h)^{d-1},
 \tag{P-F55h}
\]
where $h$ is a degree-$d$ form in $n-1$ variables.
Substitute $n=6$, $d=3$ and $h=-f$.
Since $\deg\Delta_5=80$ is even, Proposition~4.7(i)
gives $\Delta_5(-f)=\Delta_5(f)$. Thus (P-F55h)
is exactly (P-F55d) with $C=3^{33}$ in that normalization;
arbitrary nonzero normalizations change only $C$.

Finally P-F41 and P-F46 identify the restriction of this very
quasi-pullback to the cyclic-cubic-threefold ball with $\phi$,
up to scalar, and identify its invariant image with $\Delta_5$.
Squaring these identities gives (P-F55e). On the algebraic
side (P-F55d) at $a=1$ gives the same squared discriminant.
The gradings are consistent: the restricted square has weight
$96$, five-variable coefficient degree $(5/3)96=160$, and
five-variable stretching degree $160/5=32$.
These differ from the six-variable coefficient and stretching
degrees $192,96$, which is why no equal-degree full-ring
substitution assertion has been made.
\end{proof}

\paragraph{Scope for the theta comparison.}
P-F51 realizes the definite-$U(3)$ theta factor of the same
input $g$ used here, and P-F53 models its coefficient dimensions.
P-F55 identifies another actual invariant output and the exact
cyclic restriction of that output. The operation still includes
division by $\delta$, regularized multiplicative lifting, and
the prescribed quasi-pullback or restriction. It supplies no
identification of these operations with the supplied Remark's
iterated totally definite unitary lifts, and no basis or positive
dimension formula for all graded pieces of either invariant
ring. In particular it gives no growing-dimension PH2/SH bound.
